\section{A transferred asset can still contain an option exposure}
\label{sec:bank-option-transfers}

The guarantee example retains a possible credit loss. A call instead retains
the right to buy an asset back; a written put obliges the bank to buy it if
the counterparty exercises. Neither name determines whether a sale has
occurred for accounting purposes. The transfer assessment first identifies
the cash flows being tested and whether substantially all risks and rewards
were retained or transferred. Only the intermediate case proceeds to control
and, where control remains, continuing involvement. A transferee's practical
ability to sell the whole asset unilaterally, without additional restrictions,
matters. In particular, a call over an asset readily obtainable in the market
need not prevent derecognition. Deeply in-the-money and out-of-the-money
options can also lead to different risk-transfer findings
\citep[IFRS~9.3.2.2--16, B3.2.6--9, B3.2.16]{eu20231803}.

The following fair-value cases assume that this assessment requires continuing
involvement and that the underlying asset is measured through profit or loss.
They use the historical IFRS~9 text retained in EU Regulation~2023/1803.
Amounts are in one currency, with no tax, coupons or transaction costs. Fair
values and option time values are supplied measurement inputs. Applying the
cases to an actual transfer requires its rights, market access, applicable
rules and valuation evidence.

\subsection{The asset less the liability represents the retained call}

Suppose an asset worth 80 is transferred for 75 while the bank retains a call
with exercise price 95 and value 5. The call is out of the money: exercise
would currently cost more than buying the asset. Its entire value of 5 is
therefore time value. Under the admitted continuing-involvement branch,
the bank keeps an asset of 80 and recognises an associated liability of 75.
Their difference is precisely the call's value. Opening cash 200, asset 80
and equity 280 become cash 275, asset 80, liability 75 and unchanged equity.
The receipt debits cash and credits the associated liability; it creates no
sale gain in this constructed example.

Let \(S\) denote the asset's current fair value, \(K\) the call exercise
price and \(T_c\) its supplied time value, all in the same currency. The
call's intrinsic value is \(\max(S-K,0)\). The source's two measurement
branches can consequently be written together as
\[
 A=S,\qquad L=\min(S,K)-T_c,\qquad
 A-L=\max(S-K,0)+T_c.
\]
Here \(A\) and \(L\) are the recognised asset and associated liability.
When the call moves into the money, the liability uses the exercise price
instead of the higher asset value. This preserves the same net call exposure
\citep[IFRS~9.B3.2.13(c)]{eu20231803}.

At a later measurement, take \(S=100\) and \(T_c=4\). The asset becomes
100, the liability 91 and the net exposure 9. Asset income 20 and liability
expense 16 produce net profit 4. At expiry, take \(S=110\) and zero time
value. The asset becomes 110 and the liability 95, producing further net
profit \(10-4=6\). Physical exercise pays 95 and removes the liability;
the already recognised asset remains. Closing cash 180 plus asset 110 equals
equity 290, the original 280 plus profit 10. Recognising a separate call of
5 at transfer would have counted the same right twice. IFRS~9.B3.2.14
prohibits that duplication, while 3.2.22 preserves separate asset/liability
and income/expense presentation \citep[IFRS~9.3.2.19, 3.2.22, B3.2.14]{eu20231803}.

\subsection{A written put caps the retained asset; a collar combines rights}

For a written put, the bank bears the downside below the exercise price but
does not retain the upside above it. With exercise price \(K\) and put time
value \(T_p\), the source therefore limits the recognised asset and measures
the associated liability as
\[
 A=\min(S,K),\qquad L=K+T_p,\qquad
 A-L=-\bigl(\max(K-S,0)+T_p\bigr).
\]
The negative net amount is the full put obligation
\citep[IFRS~9.B3.2.13(d)]{eu20231803}. For \(S=120\), \(K=100\) and
\(T_p=5\), asset 100 and liability 105 represent a net obligation of 5.
Choose transfer proceeds 125, equal to the underlying value plus compensation
for that obligation. From cash 200, asset 120 and equity 320, the entry debits
cash 125, credits asset 20 and credits liability 105. Equity stays 320.
Keeping the full asset value 120 would overstate the retained asset by 20.

If the asset falls to 90 and time value is 4, the recognised asset is 90 and
liability 104. Net loss is \(10-1=9\). At expiry with asset value 80,
the liability becomes 100 and the further loss is \(10-4=6\). Exercise
pays 100 and removes that liability. Cash 225 plus asset 80 equals equity
305, reconciling the total loss 15. This is a physical-settlement case;
different cash-settlement rights require their own derecognition analysis.

A collar combines a retained call at an upper exercise price \(K_c\) with
a written put at a lower price \(K_p\). Let \(P\) be the \emph{full}
put value, \(\max(K_p-S,0)+T_p\). The asset continues at fair value, and
the associated liability is
\[
 A=S,\qquad L=\min(S,K_c)+P-T_c,\qquad
 A-L=\bigl(\max(S-K_c,0)+T_c\bigr)-P.
\]
Thus the recognised net position is the call less the put
\citep[IFRS~9.B3.2.13(e)]{eu20231803}. With \(S=100\), call strike 120,
put strike 80 and time values 5 and 1, the liability is 96 and net asset 4.
From cash 200, asset 100 and equity 300, proceeds 96 produce cash 296,
asset 100 and liability 96 without an initial gain.

Two subsequent paths show why the full put value matters. At \(S=130\),
call time value 2 and put time value 1 give liability 119 and net asset 11.
At \(S=70\), call time value 1 and put time value 2 give full put value
12, liability 81 and net liability 11. Using only the put's time value in
the latter calculation would understate the associated liability by 10.
If these respective asset values persist to expiry, both time values vanish.
The upper path pays 120 to exercise the call, closing with cash 176, asset
130 and equity 306. The lower path pays 80 on the exercised put, closing
with cash 216, asset 70 and equity 286. Cumulative income is respectively
6 and \(-14\), measured from the initial net option value 4.

\subsection{Amortised cost uses two interest schedules}

The fair-value formulas cannot be copied into an amortised-cost case.
IFRS~9.B3.2.13(b) instead accretes the associated liability from proceeds to
the transferred asset's gross carrying amount at option expiry, using
effective interest. Any difference between that liability and the exercise
price enters profit when the option is exercised
\citep[IFRS~9.B3.2.13(b)]{eu20231803}.

Take the source's gross asset 98, proceeds 95 and gross asset 100 at expiry,
here one year later. Assume a zero loss allowance, no intervening coupons
or cash flows, and collection of the asset after option expiry. The asset's
effective rate is \(100/98-1=1/49\), producing income 2. The liability's
rate is \(100/95-1=1/19\), producing expense 5. They cannot share a
single rate merely because they concern the same transfer. If a written put
is then exercised at 102, debit liability 100 and loss 2, and credit cash
102. Opening cash 200, asset 98 and equity 298 become cash 193, asset 100
and equity 293. Interest expense net of income is 3 and exercise loss is
2. The zero allowance isolates the financing calculation; it does not infer
an absence of credit risk from the transfer.
